%%%%%%%%%%  scan idx 86  |  folio 73  %%%%%%%%%%
Si \(G\) est un
faisceau cohérent sur \(Y\), \(f_*(\underline G)\) est cohérent sur \(X\). Si l'on note \(R^pf_*\) le \(p\)-ème foncteur dérivé du foncteur \(f\) ainsi défini, on a, par définition, \(f_!(\underline G)=\sum(-1)^p R^pf_*(\underline G)\).\quad*

\textbf{THÉORÈME DE RIEMANN--ROCH.} \emph{Soit \(f:Y\longrightarrow X\) un morphisme} (* application régulière *) \emph{propre d'une variété quasi-projective non singulière \(Y\) dans une autre \(X\), et soit \(y\in K(Y)\). On a~:}
\begin{equation}\tag{1}
  f_*\bigl(\ch(y)T(Y)\bigr)=\ch(f_!(y))T(X)
  \quad\text{\emph{dans} } A(X)\otimes\mathbb Q.
\end{equation}

\emph{Démonstration} (* abrégée *).

\textbf{LEMME 1.} \emph{Considérons des morphismes propres de variétés \(Z\xrightarrow{g}Y\xrightarrow{f}X\). Si le théorème \(\RR\) est vrai pour \(f\) et \(g\), il l'est pour \(f\circ g\). S'il est vrai pour \(fg\) et \(g\), il l'est pour \(f\). Si \(f_*\) est injective, et si le théorème \(\RR\) est vrai pour \(fg\) et \(f\), il l'est pour \(g\).}

Trivial. (* En fait, lorsqu'on suppose \(\RR\) vrai pour \(fg\) et \(g\), il faut encore supposer que \(g_!\) est surjectif pour pouvoir affirmer que \(\RR\) est vrai pour \(f\). *)

Un morphisme propre étant composé d'une injection et d'une projection \(X\times P\longrightarrow X\) (où \(P\) est un espace projectif), on est ramené à traiter séparément ces deux cas. Le deuxième, on l'a vu, se ramène à du calcul standard, moyennant la proposition 2.13 (« décomposition cellulaire ») de R.R.R. On va donc supposer \(Y\subset X\), avec le morphisme d'injection \(i\). La formule (1) devient alors~:
\begin{equation}\tag{2}
  \ch(i_!(y))=i_*\bigl(\ch(y)T(E)^{-1}\bigr),
\end{equation}
où \(E=T_{X/Y}\) est le fibré normal à \(Y\) dans \(X\).

